% Conclusion
We studied fresh data collection by a UAV swarm that recharges at a small number
of shared, multi-port charging stations. We formulated the joint placement,
capacity, trajectory, and assignment problem under a charging queue, and we
showed that the correct queue model is a finite-source (machine-repair) queue,
not an open $M/M/c$: a discrete-event simulation confirms the finite-source
formula within about $4\%$ and shows that the open queue overpredicts the wait
by roughly $6\times$ and distorts the design. Two structural results followed: a
congestion threshold that makes the AoI non-monotone in the swarm size, with a
provisioning law linking the optimal swarm size to the charging capacity, and a
placement inversion in which the AoI-optimal placement becomes traffic-driven
rather than coverage-driven. A block-coordinate solver with a GPU
close-enough-TSP trajectory optimizer realized these gains, reducing peak AoI
across field sizes against strong baselines.

Several limitations remain. The finite-source model assumes exponential
operating and service times; the deterministic-service $M/D/c//N$ approximation
we introduce tightens the bound, but a residual gap remains from the
near-deterministic operating process, which a full non-exponential treatment
would close. The trajectory block uses a close-enough-TSP heuristic with $2$-opt
reordering rather than an exact solver; it matches the exact optimum on the small
instances where enumeration is tractable, but tighter guarantees at scale are
open. Finally, extending the model to stochastic data arrivals and online
operation would move the design from planning toward real-time control.
